% Generated draft: requires mathematical and source review.
\documentclass[11pt]{article}
\usepackage[margin=25mm]{geometry}
\usepackage{fontspec}
\setmainfont{Latin Modern Roman}
\usepackage{amsmath,amssymb,mathtools,unicode-math}
\setmathfont{Latin Modern Math}
\usepackage{enumitem,needspace,xurl,hyperref}
\setlength{\parindent}{0pt}
\setlength{\parskip}{5pt}
\setlength{\emergencystretch}{3em}
\newcommand{\N}{\mathbb{N}}
\newcommand{\Z}{\mathbb{Z}}
\newcommand{\Q}{\mathbb{Q}}
\newcommand{\R}{\mathbb{R}}
\newcommand{\C}{\mathbb{C}}
\newcommand{\F}{\mathbb{F}}
\newcommand{\A}{\mathbb{A}}
\newcommand{\PP}{\mathbb{P}}
\newcommand{\E}{\mathbb{E}}
\newcommand{\eps}{\varepsilon}
\newcommand{\dd}{\,\mathrm{d}}
\begin{document}
\section*{UnsolvedMath ID: 30002657}
\textbf{Source problem code:} catalogue entry
\subsection{Definitions and mathematical statement}
Let \(X\) be a projective hyperkähler manifold with \(\dim X=2n\). The usual setting is over \(\C\). Under the irreducible holomorphic symplectic convention, this means that \(X\) is a smooth connected simply connected complex projective variety and \(H^0(X,\Omega_X^2)=\C\sigma\), where \(\sigma\) is everywhere nondegenerate. The supplied record does not explicitly specify the ground field or whether this particular convention for hyperkähler is intended; these are qualifications requiring review.

Write \(CH_0(X)\) for the group of algebraic zero-cycles on \(X\) modulo rational equivalence, with integral coefficients, and \([y]\) for the class of a closed point \(y\in X\). A scheme has pure dimension \(n\) if each of its irreducible components has dimension \(n\). For an integral complex projective variety, rational connectedness means that two general points can be joined by the image of a morphism \(\PP^1\to Y\). The source uses rational connectedness for subschemes without specifying reducedness, irreducibility, or an extension of this definition to more general schemes. The questions below retain that unresolved terminology.

For every projective hyperkähler manifold \(X\) of dimension \(2n\), ask:
\begin{enumerate}
\item Does there exist a nonempty closed subscheme \(Y\subseteq X\) which is rationally connected and has pure dimension \(n\)?
\item For every pair of nonempty rationally connected closed subschemes \(Y_1,Y_2\subseteq X\) with \(\dim Y_1=\dim Y_2=n\), and every choice of closed points \(y_1\in Y_1\) and \(y_2\in Y_2\), is
\[
[y_1]=[y_2]\quad\text{in }CH_0(X)?
\]
\end{enumerate}
The second question includes \(Y_1=Y_2\). If the first question has a positive answer and the second equality holds, denote the resulting common point class by \(c_X\in CH_0(X)\); it is to be independent of both the subscheme and the point. Without existence of any eligible subscheme, the universal equality is vacuous and does not define \(c_X\). The source explicitly requires pure dimension in the first question, but states only dimension \(n\) in the second; no additional purity assumption has been inserted there.
\subsection{Short English statement}
Does every projective hyperkähler manifold of dimension \(2n\) contain a rationally connected subscheme of pure dimension \(n\), and do all points on all rationally connected \(n\)-dimensional subschemes represent one common class in \(CH_0(X)\)? The meaning of rational connectedness for the allowed subschemes and the precise hyperkähler convention require clarification.
\subsection{Sources}
\begin{itemize}
\item Supplied local UnsolvedMath registry record, permanent ID \(30002657\), title \emph{Rationally Connected Cycles on Hyperkähler Manifolds}, collection \emph{Oberwolfach Reports Open Problems}, exact source problem code \texttt{catalogue entry}. No individual proposer is supplied; the proposed year is recorded as 2014.
\item Supplied report citation: \emph{Komplexe Analysis (2015)}, Oberwolfach Reports, pp.~2177--2234, DOI \texttt{10.4171/owr/2014/39}. \url{https://doi.org/10.4171/owr/2014/39}. This is also the supplied literature source labelled ``original Oberwolfach report.'' The supplied report title year is 2015, whereas the DOI contains 2014 and the registry records a proposed year of 2014; their relationship has not been independently checked.
\item Supplied literature source: Charles and Pacienza, ``rational curves and zero-cycles.'' \url{https://arxiv.org/abs/1401.4071}. This descriptive label is retained as supplied; no additional bibliographic details or theorem attribution are inferred.
\item Literature assessment checked at \texttt{2026-08-21}: the supplied dated curation labels the problem open and reports positive constructions in special hyperkähler deformation types and dimensions, while asserting that the existence and universal \(CH_0\)-class questions remain open for arbitrary projective hyperkähler manifolds. This is supplied dated curation derived from the OWR record, not an independently verified literature assessment. The record itself requests independent verification of time-sensitive status and cited claims.
\item The registry labels the statement \texttt{corrected\_verified} and says its mathematical content was unchanged while notation, typography, and wording were normalized. These are supplied editorial assertions, not independent mathematical verification of this draft.
\end{itemize}
\end{document}
